% A1_brief.tex -- ONE source for the English and Romanian brief of Assignment 1.
% Compile A1_brief_EN.tex or A1_brief_RO.tex (they set the language and \input this file).
\renewcommand{\emqaassign}{\T{Assignment 1}{Tema 1}}
\begin{document}
\emqatitle{\T{Assignment 1: Forecasting the Romanian day-ahead price for a bidding desk}%
                {Tema 1: Prognoza prețului PZU din România, pentru un birou de ofertare}}%
          {\T{Point and probabilistic forecasts, forecast tests and volatility --- lectures 1--3}%
             {Prognoze punctuale și probabilistice, teste de acuratețe și volatilitate --- cursurile 1--3}}

\begin{todobox}\small
\textbf{\T{Instructor to fill in}{De completat de profesor}:}
\T{Deadline}{Termen}: \textbf{TODO} \quad\textbullet\quad
\T{Weight in the final grade}{Pondere în nota finală}: \textbf{TODO} \quad\textbullet\quad
\T{Individual or teams of 2}{Individual sau echipe de cîte 2}: \textbf{TODO} \quad\textbullet\quad
\T{Oral defence slots}{Programarea susținerilor orale}: \textbf{TODO} \quad\textbullet\quad
\T{Late policy}{Regula pentru întîrzieri}: \textbf{TODO}.
\T{Submit on Moodle.}{Predarea se face pe Moodle.}
\end{todobox}

\begin{scenario}
\T{You work on the day-ahead desk of a Romanian energy company: a supplier, a producer or a large consumer.
Every day, before 12:00 CET (Central European Time) on day $d-1$, the desk sends bids to the Romanian day-ahead market
(PZU, \emph{Piața pentru Ziua Următoare}, operated by OPCOM, the Romanian power and gas market operator) for each hour of day $d$.
Today the desk prices its bids with ``the price of yesterday'' (a naive rule).
Your manager asks three questions.
(1) Should we replace the naive rule by a statistical model, at least in the hours that matter most for us (your \emph{focus block})?
(2) How wide is the uncertainty we should communicate with every forecast?
(3) Is the market now in a high-volatility regime that justifies a larger credit or collateral buffer?
You answer with a reproducible notebook and a two-page memo that ends in a clear decision.}%
{Lucrați la biroul de ofertare pentru ziua următoare al unei companii energetice din România: un furnizor, un producător sau un consumator mare.
În fiecare zi, înainte de ora 12:00 CET (ora Europei Centrale) din ziua $d-1$, biroul trimite oferte pe piața pentru ziua următoare
(PZU; în engleză \emph{day-ahead market}, operată de OPCOM, operatorul pieței de energie electrică și de gaze naturale din România) pentru fiecare oră a zilei $d$.
Astăzi biroul își stabilește ofertele cu „prețul de ieri'' (o regulă naivă).
Managerul vă pune trei întrebări.
(1) Înlocuim regula naivă cu un model statistic, măcar în orele care contează cel mai mult pentru noi (\emph{blocul dumneavoastră de interes})?
(2) Cît de mare este incertitudinea pe care trebuie să o comunicăm odată cu fiecare prognoză?
(3) Este piața acum într-un regim de volatilitate ridicată, care justifică o rezervă mai mare de credit sau de garanții?
Răspundeți cu un notebook reproductibil și cu un memoriu de două pagini care se încheie cu o decizie clară.}
\end{scenario}

\section{\T{What you submit}{Ce predați}}
\begin{enumerate}[label=\textbf{D\arabic*}]
\item \T{\textbf{Notebook} \texttt{A1\_<ID>.ipynb}, built from the starter \texttt{A1\_starter\_EN.ipynb}. It must run with \emph{Restart \& Run All}
  from the course data folder, without errors, in at most 30 minutes on a normal laptop.}%
  {\textbf{Notebook-ul} \texttt{A1\_<ID>.ipynb}, construit pornind de la \texttt{A1\_starter\_RO.ipynb}. Trebuie să ruleze cu \emph{Restart \& Run All}
  din dosarul de date al cursului, fără erori, în cel mult 30 de minute pe un laptop obișnuit.}
\item \T{\textbf{Numbers file} \texttt{my\_numbers\_A1\_<ID>\_personal.json}, written by the notebook (last cell of the starter).}%
  {\textbf{Fișierul de numere} \texttt{my\_numbers\_A1\_<ID>\_personal.json}, scris de notebook (ultima celulă a notebook-ului de pornire).}
\item \T{\textbf{Memo} (PDF, 11 pt, at most 2 pages): page 1 is the recommendation for your manager; page 2 is the appendix
  (key tables and the hand-reconciliation table). Structure in Task 8.}%
  {\textbf{Memoriul} (PDF, 11 pt, cel mult 2 pagini): pagina 1 este recomandarea pentru manager; pagina 2 este anexa
  (tabelele principale și tabelul de reconciliere manuală). Structura este în sarcina 8.}
\item \T{\textbf{Oral defence} (about 10 minutes): the instructor picks one day and one hour of your test window and one adverse question;
  you explain your numbers from your own notebook.}%
  {\textbf{Susținerea orală} (circa 10 minute): profesorul alege o zi și o oră din fereastra dumneavoastră de test și o întrebare dificilă;
  explicați numerele din propriul notebook.}
\end{enumerate}

\section{\T{Your personal parameters}{Parametrii dumneavoastră personali}}
\T{Every student (or team) receives different settings. They are computed from your student ID by a fixed rule (a seed from the SHA-256 secure hash algorithm),
so you always get the same values and nobody else gets the same set. Get them in either way:}%
{Fiecare student (sau echipă) primește alți parametri. Ei sînt calculați din numărul matricol printr-o regulă fixă (o valoare inițială, \emph{seed}, obținută cu algoritmul de hash SHA-256, \emph{Secure Hash Algorithm}),
deci obțineți mereu aceleași valori și nimeni altcineva nu are același set. Îi obțineți în oricare dintre aceste două moduri:}
\begin{itemize}
\item \texttt{python make\_assignment.py -{}-id <ID> -{}-assignment 1} \T{(folder \texttt{repo/assignments}); team: \texttt{-{}-id "ID1+ID2"};}{(dosarul \texttt{repo/assignments}); echipă: \texttt{-{}-id "ID1+ID2"};}
\item \T{or the first cell of the starter notebook (\texttt{STUDENT\_ID = ...}).}{sau prima celulă a notebook-ului de pornire (\texttt{STUDENT\_ID = ...}).}
\end{itemize}
\begin{center}\small
\begin{tabularx}{\linewidth}{>{\ttfamily}l X}
\toprule
\T{\normalfont\textbf{Parameter}}{\normalfont\textbf{Parametru}} & \T{\textbf{Meaning}}{\textbf{Semnificație}}\\
\midrule
test\_start, test\_end & \T{your 365-day test window (first day of a month between Sep 2024 and Nov 2025)}{fereastra dumneavoastră de test, de 365 de zile (începe în prima zi a unei luni între sept.\ 2024 și nov.\ 2025)}\\
burn\_in\_days & \T{112 days before the test window that you also forecast; their errors feed the probabilistic forecasts (Task 6)}{112 zile dinaintea ferestrei de test, pe care le prognozați de asemenea; erorile lor alimentează prognozele probabilistice (sarcina 6)}\\
focus\_hours & \T{the hours that matter for your desk (CET): night, morning ramp, solar hours or evening peak}{orele care contează pentru biroul dumneavoastră (CET): noapte, rampa de dimineață, orele solare sau vîrful de seară}\\
arx\_windows & \T{the six calibration windows 56, 84, 112, 714, 721, 728 days (the same for everybody)}{cele șase ferestre de calibrare: 56, 84, 112, 714, 721, 728 de zile (aceleași pentru toți)}\\
lear\_window & \T{your calibration window for LEAR: 364, 546 or 728 days}{fereastra dumneavoastră de calibrare pentru LEAR: 364, 546 sau 728 de zile}\\
error\_window & \T{$W$: number of past days of errors used for HS and CP intervals: 56, 84 or 112}{$W$: numărul de zile anterioare ale căror erori se folosesc pentru intervalele HS și CP: 56, 84 sau 112}\\
garch\_years & \T{length of the GARCH estimation sample before test\_start: 3, 4 or 5 years}{lungimea eșantionului de estimare GARCH dinainte de test\_start: 3, 4 sau 5 ani}\\
check\_days, check\_hour & \T{two days and one hour you must reconcile by hand (Task 8)}{două zile și o oră pe care trebuie să le reconciliați manual (sarcina 8)}\\
\bottomrule
\end{tabularx}
\end{center}
\T{Copy the printed block, unchanged, into the memo appendix. A memo with someone else's parameters receives 0 points.}%
{Copiați blocul afișat, nemodificat, în anexa memoriului. Un memoriu cu parametrii altcuiva primește 0 puncte.}

\section{\T{Data, time and notation}{Date, timp și notație}}
\T{All data are in \texttt{repo/data/clean/} (sources and download dates in \texttt{repo/data/SOURCES.md}).
Use the helper \texttt{emqa\_assign.py}; do not download other data and do not edit the data files.}%
{Toate datele sînt în \texttt{repo/data/clean/} (sursele și data descărcării în \texttt{repo/data/SOURCES.md}).
Folosiți modulul \texttt{emqa\_assign.py}; nu descărcați alte date și nu modificați fișierele de date.}
\begin{center}\small
\begin{tabularx}{\linewidth}{l l X}
\toprule
\T{\textbf{Object}}{\textbf{Obiect}} & \T{\textbf{Call}}{\textbf{Apel}} & \T{\textbf{Content}}{\textbf{Conținut}}\\
\midrule
$P_{d,h}$ & \texttt{ea.panel("RO")[0]} & \T{day-ahead price of Romania, day $d$, hour $h$ (EUR/MWh)}{prețul PZU al României, ziua $d$, ora $h$ (EUR/MWh)}\\
$L_{d,h}$ & \texttt{ea.panel("RO")[1]} & \T{actual load of Romania (mean MW in the hour)}{consumul realizat al României (MW medii în oră)}\\
$R_{d,h}$ & \texttt{ea.panel("RO")[2]} & \T{actual wind + solar generation of Romania (MW)}{producția eoliană + solară realizată a României (MW)}\\
$\bar P_d$ & \texttt{ea.daily\_prices()["RO"]} & \T{daily base price: mean of all market time units of the CET day}{prețul zilnic de bază: media tuturor intervalelor de tranzacționare ale zilei CET}\\
\bottomrule
\end{tabularx}
\end{center}
\begin{itemize}
\item \T{\textbf{Days and hours.} One row per delivery day $d$ in CET (Central European Time, the time zone of the European market coupling, SDAC);
columns \texttt{h00}--\texttt{h23} are CET hours, $h=0,\dots,23$. Romanian local time (EET, Eastern European Time) is one hour later: \texttt{h00} is 01:00--02:00 in Bucharest.}%
{\textbf{Zile și ore.} Un rînd pe zi de livrare $d$ în ora CET (ora Europei Centrale, fusul orar al cuplării unice a piețelor pentru ziua următoare, SDAC);
coloanele \texttt{h00}--\texttt{h23} sînt ore CET, $h=0,\dots,23$. Ora României (EET, ora Europei de Est) este cu o oră mai tîrziu: \texttt{h00} înseamnă 01:00--02:00 la București.}
\item \T{\textbf{15-minute market time unit (MTU).} Since delivery day 1 October 2025 the market clears 96 quarter-hours per day.
In the hourly matrix each hour after that date is the mean of its four quarter-hours. Forecast and evaluate hourly values.}%
{\textbf{Intervalul de tranzacționare de 15 minute (MTU, \emph{market time unit}).} Din ziua de livrare 1 octombrie 2025 prețurile se stabilesc pentru 96 de sferturi de oră pe zi.
În matricea orară, fiecare oră de după această dată este media celor patru sferturi de oră. Prognozați și evaluați valori orare.}
\item \T{\textbf{Information set (gate closure).} A forecast of day $d$ is made at 12:00 CET on day $d-1$. At that time:
\begin{itemize}\setlength{\itemsep}{0pt}
\item \emph{prices}: all prices up to day $d-1$ are known (the $d-1$ auction cleared on day $d-2$);
\item \emph{realised load $L$ and wind+solar $R$}: only values already metered, i.e.\ up to the end of day $d-2$; day $d-1$ is \emph{not} yet complete;
\item \emph{fuels and carbon}: the last close before 12:00 on $d-1$, in practice the close of day $d-2$;
\item \emph{transformations, scalers, thresholds}: fitted on these data only.
\end{itemize}
A real desk would use the forecasts of the transmission system operator (TSO, Transelectrica); our free data contain only actual values.
\textbf{Rule:} use the actual $L$ and $R$ of day $d-2$ (same hour) as the proxy for day $d$, and $d-3$, $d-7$ as further lags. \texttt{ea.leakage\_check} enforces this rule.}%
{\textbf{Setul de informații (închiderea porții, \emph{gate closure}).} Prognoza pentru ziua $d$ se face la 12:00 CET în ziua $d-1$. În acel moment:
\begin{itemize}\setlength{\itemsep}{0pt}
\item \emph{prețurile}: toate prețurile pînă în ziua $d-1$ sînt cunoscute (licitația pentru ziua $d-1$ s-a încheiat în ziua $d-2$);
\item \emph{consumul $L$ și producția eoliană+solară $R$ realizate}: numai valorile deja măsurate, adică pînă la sfîrșitul zilei $d-2$; ziua $d-1$ \emph{nu} este încă încheiată;
\item \emph{combustibilii și certificatele de emisii}: ultima cotație de închidere dinainte de ora 12:00 din ziua $d-1$, în practică cea din ziua $d-2$;
\item \emph{transformările, standardizările și pragurile}: estimate numai pe aceste date.
\end{itemize}
Un birou real ar folosi prognozele operatorului de transport și sistem (OTS, Transelectrica); datele noastre gratuite conțin doar valori realizate.
\textbf{Regulă:} folosiți $L$ și $R$ realizate din ziua $d-2$ (aceeași oră) ca substitut (proxy) pentru ziua $d$, iar $d-3$ și $d-7$ ca decalaje suplimentare. \texttt{ea.leakage\_check} verifică respectarea acestei reguli.}
\item \T{\textbf{Daylight saving time (DST).} The CET day has 23 hours in March and 25 in October. The helper keeps 24 columns:
the missing hour is the mean of its neighbours, the doubled hour is averaged (column \texttt{dst} marks these days).}%
{\textbf{Ora de vară (DST).} Ziua CET are 23 de ore în martie și 25 în octombrie. Modulul păstrează 24 de coloane:
ora lipsă este media orelor vecine, ora dublată este mediată (coloana \texttt{dst} marchează aceste zile).}
\end{itemize}

\section{\T{Tasks}{Sarcinile}}
\T{Do the tasks in order. Each task says what to compute, on which data, with which method, and what to show.
The points add up to 100.}%
{Rezolvați sarcinile în ordine. Fiecare sarcină spune ce calculați, pe ce date, cu ce metodă și ce arătați.
Punctele însumează 100.}

\task{1}{\T{Know your test window}{Cunoașteți-vă fereastra de test}}{8}
\T{On the 365 test days, compute the mean, standard deviation, minimum and maximum of $P_{d,h}$; the number of negative hours;
the number of negative market time units (MTUs) in \texttt{ea.quarter\_hours("RO")} (hourly rows before 1 October 2025, quarter-hour rows after; filter on the CET date of the index); and the share of test days with 15-minute MTU.
Plot the daily base price from 1 January 2023 to the end of your test window and shade the window.
In three sentences, name the market events inside your window (for example the Romania--Hungary (RO--HU) price separation of summer 2024,
the dry hydro year 2025, negative prices in spring and summer, the 15-minute MTU from October 2025).}%
{Pe cele 365 de zile de test calculați media, abaterea standard, minimul și maximul lui $P_{d,h}$; numărul de ore cu preț negativ;
numărul intervalelor de tranzacționare (MTU) cu preț negativ din \texttt{ea.quarter\_hours("RO")} (rînduri orare înainte de 1 octombrie 2025, sferturi de oră după; filtrați după data CET a indexului); ponderea zilelor de test cu MTU de 15 minute.
Reprezentați grafic prețul zilnic de bază de la 1 ianuarie 2023 pînă la sfîrșitul ferestrei de test și hașurați fereastra.
În trei propoziții, numiți evenimentele de piață din fereastra dumneavoastră (de exemplu separarea prețurilor România--Ungaria (RO--HU) din vara 2024,
anul hidrologic secetos 2025, prețurile negative din primăvară și vară, MTU de 15 minute din octombrie 2025).}
\out{\T{table T1 (one column); figure with x-axis ``date'', y-axis ``EUR/MWh'', legend below the plot.}{tabelul T1 (o coloană); figura cu axa x „data'', axa y „EUR/MWh'', legenda sub grafic.}}

\task{2}{\T{Two naive benchmarks}{Două modele naive de referință}}{6}
\T{For every day of burn-in + test (\texttt{fdays}) and every hour:}%
{Pentru fiecare zi din perioada de rodaj (\emph{burn-in}) + test (\texttt{fdays}) și fiecare oră:}
\begin{equation}
\hat P^{\text{naive}}_{d,h}=\begin{cases}P_{d-7,h}, & d=\text{\T{Mon, Sat, Sun}{luni, sîmbătă, duminică}}\\ P_{d-1,h}, & d=\text{\T{Tue--Fri}{marți--vineri}}\end{cases}
\qquad\qquad \hat P^{(7)}_{d,h}=P_{d-7,h}.
\end{equation}
\T{The first is the naive benchmark of the electricity-price-forecasting (EPF) literature \cite{Weron2014,Lago2021}; the second is the seasonal naive.
Every later model is judged against $\hat P^{\text{naive}}$.}%
{Primul este modelul naiv de referință din literatura de prognoză a prețului energiei electrice (EPF, \emph{electricity price forecasting}) \cite{Weron2014,Lago2021}; al doilea este modelul naiv sezonier.
Toate modelele următoare sînt evaluate în raport cu $\hat P^{\text{naive}}$.}
\out{\T{both forecasts stored in \texttt{F} (\texttt{DataFrame}, rows \texttt{fdays}, columns \texttt{h00}..\texttt{h23}); \texttt{ea.leakage\_check} passed.}{ambele prognoze salvate în \texttt{F} (\texttt{DataFrame}, rînduri \texttt{fdays}, coloane \texttt{h00}..\texttt{h23}); testul \texttt{ea.leakage\_check} reușit.}}

\task{3}{\T{Expert ARX model, rolling windows and window averaging}{Modelul expert ARX, ferestre glisante și medierea pe ferestre}}{12}
\T{\textbf{Transform.} Prices have spikes and negative values, so never take logarithms of prices.
Use the variance-stabilising transformation (VST) of \cite{Uniejewski2018}:}%
{\textbf{Transformarea.} Prețurile au salturi și valori negative, deci nu logaritmați niciodată prețurile.
Folosiți transformarea de stabilizare a varianței (VST, \emph{variance-stabilising transformation}) din \cite{Uniejewski2018}:}
\begin{equation}
Y_{d,h}=\asinh\!\left(\frac{P_{d,h}-m}{s}\right),\qquad \asinh(x)=\ln\!\big(x+\sqrt{x^2+1}\big),\qquad P_{d,h}=m+s\,\sinh(\hat Y_{d,h}),
\end{equation}
\T{\where\ $m$ is the median and $s$ the median absolute deviation (MAD, scaled by 1\ds 4826) of the prices in the calibration window only.}%
{\where\ $m$ este mediana și $s$ abaterea absolută mediană (MAD, înmulțită cu 1\ds 4826) a prețurilor numai din fereastra de calibrare.}

\T{\textbf{Model.} For each hour $h$ separately (24 regressions per day) estimate by ordinary least squares (OLS) the expert ARX model
(autoregressive model with exogenous variables) of \cite{Ziel2018}:}%
{\textbf{Modelul.} Pentru fiecare oră $h$ separat (24 de regresii pe zi) estimați prin metoda celor mai mici pătrate (OLS, \emph{ordinary least squares}) modelul expert ARX
(model autoregresiv cu variabile exogene) din \cite{Ziel2018}:}
\begin{multline}
Y_{d,h}=\beta_{h,0}+\beta_{h,1}Y_{d-1,h}+\beta_{h,2}Y_{d-2,h}+\beta_{h,3}Y_{d-7,h}+\beta_{h,4}Y_{d-1,23}+\beta_{h,5}Y^{\min}_{d-1}+\beta_{h,6}Y^{\max}_{d-1}\\
+\beta_{h,7}\tilde L_{d,h}+\beta_{h,8}\tilde R_{d,h}+\beta_{h,9}D^{\text{Mon}}_d+\beta_{h,10}D^{\text{Sat}}_d+\beta_{h,11}D^{\text{Sun}}_d+\varepsilon_{d,h}
\end{multline}
\T{\where\ $Y_{d-1,23}$ is the last hour of the previous day; $Y^{\min}_{d-1},Y^{\max}_{d-1}$ are the minimum and maximum over the 24 hours of day $d-1$;
$\tilde L_{d,h}=L_{d-2,h}$ and $\tilde R_{d,h}=R_{d-2,h}$ are the proxies (the last metered day; standardised with the mean and standard deviation of the calibration window, rows up to $d-2$);
$D$ are 0/1 dummies for Monday, Saturday and Sunday; $\varepsilon_{d,h}$ is the error.}%
{\where\ $Y_{d-1,23}$ este ultima oră a zilei precedente; $Y^{\min}_{d-1},Y^{\max}_{d-1}$ sînt minimul și maximul celor 24 de ore din ziua $d-1$;
$\tilde L_{d,h}=L_{d-2,h}$ și $\tilde R_{d,h}=R_{d-2,h}$ sînt substitutele (ultima zi măsurată; standardizate cu media și abaterea standard din fereastra de calibrare, rîndurile pînă în $d-2$);
$D$ sînt variabile 0/1 pentru luni, sîmbătă și duminică; $\varepsilon_{d,h}$ este eroarea.}

\T{\textbf{Rolling calibration.} For each forecast day $d$ and each window length $T\in\{56,84,112,714,721,728\}$, estimate on the $T$ target days
$d-T,\dots,d-1$, then forecast day $d$; move one day forward and re-estimate (daily recalibration).
\textbf{Averaging.} The combined forecast is the simple mean over the six windows,
$\hat P^{\text{AW}}_{d,h}=\tfrac16\sum_{T}\hat P^{(T)}_{d,h}$ \cite{Hubicka2019}.}%
{\textbf{Calibrarea glisantă.} Pentru fiecare zi de prognoză $d$ și fiecare lungime de fereastră $T\in\{56,84,112,714,721,728\}$, estimați pe cele $T$ zile-țintă
$d-T,\dots,d-1$, apoi prognozați ziua $d$; avansați o zi și reestimați (recalibrare zilnică).
\textbf{Medierea.} Prognoza combinată este media simplă pe cele șase ferestre,
$\hat P^{\text{AW}}_{d,h}=\tfrac16\sum_{T}\hat P^{(T)}_{d,h}$ \cite{Hubicka2019}.}
\out{\T{seven forecasts in \texttt{F}: ARX(56), \dots, ARX(728) and AW-ARX; \texttt{ea.leakage\_check} passed for ARX(56).}{șapte prognoze în \texttt{F}: ARX(56), \dots, ARX(728) și AW-ARX; testul \texttt{ea.leakage\_check} reușit pentru ARX(56).}}

\task{4}{\T{LEAR, the strong benchmark}{LEAR, modelul de referință puternic}}{10}
\T{LEAR (LASSO-estimated autoregression; LASSO = least absolute shrinkage and selection operator) \cite{Uniejewski2016,Lago2021}
uses many candidate regressors and lets the LASSO penalty switch off the useless ones. For each hour $h$:}%
{LEAR (autoregresie estimată cu LASSO; LASSO = \emph{least absolute shrinkage and selection operator}) \cite{Uniejewski2016,Lago2021}
folosește mulți regresori candidați și lasă penalizarea LASSO să-i anuleze pe cei inutili. Pentru fiecare oră $h$:}
\begin{equation}
\hat{\bm\beta}_h=\arg\min_{\bm\beta}\Big\{\sum_{d\in\text{cal}}\big(Y_{d,h}-\mathbf x_d^\top\bm\beta\big)^2+\lambda_h\sum_{j}|\beta_j|\Big\},
\end{equation}
\T{\where\ $\mathbf x_d$ has 247 entries: $Y_{d-k,0},\dots,Y_{d-k,23}$ for $k=1,2,3,7$ (96); load and wind+solar proxies of $d-2,d-3,d-7$ for all 24 hours,
after the same $\asinh$ transform (144); 7 day-of-week dummies. All columns are standardised with calibration-window means and standard deviations.
$\lambda_h\ge 0$ is chosen by the Akaike information criterion (AIC): \texttt{LassoLarsIC(criterion="aic")} in \texttt{scikit-learn}.
Use your \texttt{lear\_window}; recalibrate daily (weekly recalibration is allowed if you say so and report the run time).}%
{\where\ $\mathbf x_d$ are 247 de elemente: $Y_{d-k,0},\dots,Y_{d-k,23}$ pentru $k=1,2,3,7$ (96); substitutele de consum și eolian+solar din $d-2,d-3,d-7$ pentru toate cele 24 de ore,
după aceeași transformare $\asinh$ (144); 7 variabile pentru ziua săptămînii. Toate coloanele se standardizează cu mediile și abaterile standard din fereastra de calibrare.
$\lambda_h\ge 0$ se alege cu criteriul informațional Akaike (AIC): \texttt{LassoLarsIC(criterion="aic")} din \texttt{scikit-learn}.
Folosiți \texttt{lear\_window}; recalibrați zilnic (recalibrarea săptămînală este permisă dacă o declarați și raportați timpul de rulare).}
\out{\T{LEAR forecasts in \texttt{F}; the number of non-zero coefficients for hour \texttt{check\_hour} on your first test day.}{prognozele LEAR în \texttt{F}; numărul de coeficienți nenuli pentru ora \texttt{check\_hour} în prima zi de test.}}

\task{5}{\T{Accuracy and Diebold--Mariano tests}{Acuratețe și testele Diebold--Mariano}}{16}
\T{On the test window only (never on the burn-in), with errors $e_{d,h}=P_{d,h}-\hat P_{d,h}$ and $N$ the number of (day, hour) pairs:}%
{Numai pe fereastra de test (niciodată pe rodaj), cu erorile $e_{d,h}=P_{d,h}-\hat P_{d,h}$ și $N$ numărul de perechi (zi, oră):}
\begin{equation}
\text{MAE}=\frac1N\sum|e_{d,h}|,\qquad \text{RMSE}=\sqrt{\frac1N\sum e_{d,h}^2},\qquad \text{rMAE}=\frac{\text{MAE}_{\text{model}}}{\text{MAE}_{\text{naive}}}.
\end{equation}
\T{MAE is the mean absolute error, RMSE the root mean squared error, rMAE the relative MAE \cite{Lago2021} (below 1: better than naive).
Do not use the mean absolute percentage error (MAPE): it explodes when prices are near zero.
Report all models: (a) all 24 hours, (b) your focus block, (c) each calendar year inside your window, (d) each quarter for Naive, AW-ARX and LEAR.}%
{MAE este eroarea absolută medie, RMSE rădăcina erorii pătratice medii, rMAE eroarea absolută medie relativă \cite{Lago2021} (sub 1: mai bun decît modelul naiv).
Nu folosiți eroarea procentuală absolută medie (MAPE): crește necontrolat cînd prețurile sînt aproape de zero.
Raportați toate modelele: (a) toate cele 24 de ore, (b) blocul de interes, (c) fiecare an calendaristic din fereastră, (d) fiecare trimestru pentru Naive, AW-ARX și LEAR.}

\T{\textbf{Diebold--Mariano (DM) test} \cite{Diebold1995}, multivariate version \cite{Ziel2018}. For models A and B and each test day $d$:}%
{\textbf{Testul Diebold--Mariano (DM)} \cite{Diebold1995}, varianta multivariată \cite{Ziel2018}. Pentru modelele A și B și fiecare zi de test $d$:}
\begin{equation}
\Delta_d=\frac1{|H|}\sum_{h\in H}|e^{A}_{d,h}|-\frac1{|H|}\sum_{h\in H}|e^{B}_{d,h}|,\qquad
\text{DM}=\frac{\bar\Delta}{\sqrt{\hat\sigma^2_{\text{NW}}/n}},\qquad p=1-\Phi(\text{DM}),
\end{equation}
\T{\where\ $H$ is the set of hours (all 24, or your focus block), $n$ the number of test days, $\bar\Delta$ the mean of $\Delta_d$,
$\hat\sigma^2_{\text{NW}}$ the Newey--West (heteroskedasticity- and autocorrelation-consistent, HAC) variance of $\Delta_d$ with 7 lags, and $\Phi$ the standard normal distribution function.
A small $p$ (e.g.\ below 0\ds 05) means: B is significantly more accurate than A. Show two $6\times6$ matrices of $p$-values (all hours; focus block)
for Naive, seasonal naive, ARX(56), ARX(728), AW-ARX, LEAR.}%
{\where\ $H$ este mulțimea orelor (toate 24 sau blocul de interes), $n$ numărul de zile de test, $\bar\Delta$ media lui $\Delta_d$,
$\hat\sigma^2_{\text{NW}}$ varianța Newey--West (robustă la heteroscedasticitate și autocorelare, HAC) a lui $\Delta_d$ cu 7 decalaje, iar $\Phi$ funcția de repartiție normală standard.
Un $p$ mic (de exemplu sub 0\ds 05) înseamnă: B este semnificativ mai precis decît A. Arătați două matrice $6\times6$ de valori $p$ (toate orele; blocul de interes)
pentru Naive, naiv sezonier, ARX(56), ARX(728), AW-ARX, LEAR.}
\out{\T{table T5 (rows = models; columns MAE, RMSE, rMAE overall, focus block, per year); quarterly MAE table; two DM matrices; two sentences: which model wins, and is the win significant?}%
{tabelul T5 (rînduri = modele; coloane MAE, RMSE, rMAE total, bloc de interes, pe ani); tabelul MAE pe trimestre; cele două matrice DM; două propoziții: care model cîștigă și este cîștigul semnificativ?}}

\task{6}{\T{Probabilistic forecasts: historical simulation and conformal prediction}{Prognoze probabilistice: simulare istorică și predicție conformală}}{14}
\T{Turn a point forecast into 99 quantiles $\hat q^{\tau}_{d,h}$, $\tau=0\ds01,\dots,0\ds99$, using the errors of the same hour on the previous $W$ days
(\texttt{error\_window}), $e_{d-W,h},\dots,e_{d-1,h}$ \cite{Nowotarski2018,Lipiecki2024}:}%
{Transformați o prognoză punctuală în 99 de cuantile $\hat q^{\tau}_{d,h}$, $\tau=0\ds01,\dots,0\ds99$, folosind erorile aceleiași ore din cele $W$ zile anterioare
(\texttt{error\_window}), $e_{d-W,h},\dots,e_{d-1,h}$ \cite{Nowotarski2018,Lipiecki2024}:}
\begin{align}
\text{HS:}&\quad \hat q^{\tau}_{d,h}=\hat P_{d,h}+Q_\tau\big(e_{d-W,h},\dots,e_{d-1,h}\big),\\
\text{CP:}&\quad \hat q^{\tau}_{d,h}=\hat P_{d,h}+\operatorname{sign}(\tau-0\ds5)\;Q_{|2\tau-1|}\big(|e_{d-W,h}|,\dots,|e_{d-1,h}|\big),
\end{align}
\T{\where\ $Q_a(\cdot)$ is the empirical $a$-quantile (\texttt{numpy.quantile}, default linear interpolation). HS (historical simulation) uses the signed errors, so the interval can be asymmetric;
CP (conformal prediction) uses absolute errors, so it is symmetric around $\hat P_{d,h}$. Apply HS and CP to Naive, AW-ARX and LEAR (6 combinations).
\textbf{Pinball score} for quantile level $\tau$ and actual price $y$, and its average, the average pinball score (APS):}%
{\where\ $Q_a(\cdot)$ este cuantila empirică de ordin $a$ (\texttt{numpy.quantile}, interpolare liniară implicită). HS (simularea istorică, \emph{historical simulation}) folosește erorile cu semn, deci intervalul poate fi asimetric;
CP (predicția conformală, \emph{conformal prediction}) folosește erorile în valoare absolută, deci intervalul este simetric în jurul lui $\hat P_{d,h}$. Aplicați HS și CP pentru Naive, AW-ARX și LEAR (6 combinații).
\textbf{Scorul pinball} pentru nivelul $\tau$ și prețul realizat $y$, și media lui, scorul pinball mediu (APS, \emph{average pinball score}):}
\begin{equation}
\text{PS}(\hat q^\tau,y,\tau)=\big(\1\{y<\hat q^\tau\}-\tau\big)\big(\hat q^\tau-y\big),\qquad
\text{APS}=\frac{1}{99\,N}\sum_{d,h}\sum_{\tau}\text{PS}(\hat q^\tau_{d,h},P_{d,h},\tau),\qquad \text{CRPS}\approx 2\cdot\text{APS},
\end{equation}
\T{\where\ $\1\{\cdot\}$ is 1 if the condition holds and 0 otherwise; CRPS is the continuous ranked probability score \cite{Gneiting2007}. Lower is better.
\textbf{Coverage}: share of (day, hour) pairs with $P_{d,h}$ inside the 50\% interval $[\hat q^{0\ds25},\hat q^{0\ds75}]$ and the 90\% interval $[\hat q^{0\ds05},\hat q^{0\ds95}]$.
\textbf{Kupiec test} \cite{Kupiec1995} on the focus block for the 90\% interval, with $x$ misses out of $n$ cases and nominal miss rate $p_0=0\ds10$:}%
{\where\ $\1\{\cdot\}$ este 1 dacă este adevărată condiția și 0 altfel; CRPS este scorul probabilistic continuu ordonat (\emph{continuous ranked probability score}) \cite{Gneiting2007}. Mai mic înseamnă mai bine.
\textbf{Acoperirea}: ponderea perechilor (zi, oră) cu $P_{d,h}$ în intervalul de 50\% $[\hat q^{0\ds25},\hat q^{0\ds75}]$ și în cel de 90\% $[\hat q^{0\ds05},\hat q^{0\ds95}]$.
\textbf{Testul Kupiec} \cite{Kupiec1995} pe blocul de interes, pentru intervalul de 90\%, cu $x$ ratări din $n$ cazuri și rata nominală de ratare $p_0=0\ds10$:}
\begin{equation}
\text{LR}_{\text{UC}}=-2\ln\frac{(1-p_0)^{n-x}p_0^{\,x}}{(1-\hat p)^{n-x}\hat p^{\,x}},\qquad \hat p=x/n,\qquad \text{LR}_{\text{UC}}\sim\chi^2_1\ \text{\T{under}{sub}}\ H_0.
\end{equation}
\out{\T{table T6 (rows = 6 combinations; columns APS, APS on focus block, coverage 50\%, coverage 90\%, coverage 90\% on focus block, Kupiec $p$-value, mean width of the 90\% interval on the focus block);
one figure: one test week, actual price and the 50\% and 90\% bands of your best combination.}%
{tabelul T6 (rînduri = 6 combinații; coloane APS, APS pe bloc, acoperire 50\%, acoperire 90\%, acoperire 90\% pe bloc, valoarea $p$ Kupiec, lățimea medie a intervalului de 90\% pe bloc);
o figură: o săptămînă de test, prețul realizat și benzile de 50\% și 90\% ale celei mai bune combinații.}}

\task{7}{\T{Volatility: GARCH, GJR-GARCH, EGARCH}{Volatilitate: GARCH, GJR-GARCH, EGARCH}}{12}
\T{GARCH is the generalised autoregressive conditional heteroskedasticity model; GJR-GARCH its asymmetric version of Glosten, Jagannathan and Runkle; EGARCH the exponential GARCH. Use daily base prices. Because prices can be zero or negative, do not use log-returns; use}%
{GARCH este modelul autoregresiv generalizat cu heteroscedasticitate condiționată (\emph{generalised autoregressive conditional heteroskedasticity}); GJR-GARCH este varianta asimetrică Glosten--Jagannathan--Runkle; EGARCH este GARCH exponențial. Folosiți prețurile zilnice de bază. Deoarece prețurile pot fi zero sau negative, nu folosiți randamente logaritmice; folosiți}
\begin{equation}
r_t=100\,\big[\asinh(\bar P_t)-\asinh(\bar P_{t-1})\big],\qquad \varepsilon_t=r_t-\mu-\delta_{\text{dow}(t)}=\sigma_t z_t,\quad z_t\sim t_\nu\ \text{\T{(standardised)}{(standardizată)}},
\end{equation}
\T{\where\ $\delta_{\text{dow}(t)}$ is the mean of $r$ on that weekday, computed on the estimation sample only (subtract it before fitting), $\sigma_t$ the conditional volatility and $\nu$ the degrees of freedom. Estimate by maximum likelihood (package \texttt{arch}) on the \texttt{garch\_years} before \texttt{test\_start}:}%
{\where\ $\delta_{\text{dow}(t)}$ este media lui $r$ în acea zi a săptămînii, calculată numai pe eșantionul de estimare (o scădeți înainte de estimare), $\sigma_t$ este volatilitatea condiționată, iar $\nu$ numărul gradelor de libertate. Estimați prin metoda verosimilității maxime (pachetul \texttt{arch}) pe cei \texttt{garch\_years} ani dinaintea lui \texttt{test\_start}:}
\begin{align}
\text{GARCH(1,1):}&\quad \sigma_t^2=\omega+\alpha\varepsilon_{t-1}^2+\beta\sigma_{t-1}^2,\\
\text{GJR(1,1,1):}&\quad \sigma_t^2=\omega+\alpha\varepsilon_{t-1}^2+\gamma\varepsilon_{t-1}^2\1\{\varepsilon_{t-1}<0\}+\beta\sigma_{t-1}^2,\\
\text{EGARCH(1,1,1):}&\quad \ln\sigma_t^2=\omega+\alpha\big(|z_{t-1}|-\E|z_{t-1}|\big)+\gamma z_{t-1}+\beta\ln\sigma_{t-1}^2.
\end{align}
\T{In GJR, $\gamma>0$ means falls raise volatility more (``leverage'', usual for stocks); $\gamma<0$ means price jumps raise volatility more (``inverse leverage'', often reported for gas and power).
In EGARCH the signs are reversed in meaning: $\gamma<0$ = falls raise volatility more. \textbf{Test, do not assume}: report $\gamma$ and its $p$-value.
Persistence is $\alpha+\beta+\gamma/2$ (GJR), $\alpha+\beta$ (GARCH), $\beta$ (EGARCH); half-life $=\ln 0\ds5/\ln(\text{persistence})$ days.
Then, with the parameters fixed, filter through the test window and compare the one-day-ahead variance forecasts $\hat\sigma^2_t$ with the proxy $\varepsilon_t^2$ by the quasi-likelihood loss
$\text{QLIKE}=\frac1n\sum_t\big(\ln\hat\sigma^2_t+\varepsilon_t^2/\hat\sigma^2_t\big)$ (lower is better) \cite{Patton2011}.
Finally translate into money the one-day-ahead forecast $\hat\sigma_T$ for the last test day $T$: the approximate 99\% one-day upward move of the base price, $\bar P_{T-1}\,[\exp(z_{0\ds99}\hat\sigma_T/100)-1]$, with $\bar P_{T-1}$ the actual base price of the day before and $z_{0\ds99}=t^{-1}_\nu(0\ds99)\sqrt{(\nu-2)/\nu}$ the 99\% quantile of the standardised $t_\nu$.}%
{În GJR, $\gamma>0$ înseamnă că scăderile cresc mai mult volatilitatea („efect de levier'', obișnuit la acțiuni); $\gamma<0$ înseamnă că salturile de preț cresc mai mult volatilitatea („efect de levier invers'', des raportat pentru gaze și energie electrică).
În EGARCH semnificația semnelor este inversă: $\gamma<0$ = scăderile cresc mai mult volatilitatea. \textbf{Testați, nu presupuneți}: raportați $\gamma$ și valoarea $p$.
Persistența este $\alpha+\beta+\gamma/2$ (GJR), $\alpha+\beta$ (GARCH), $\beta$ (EGARCH); timpul de înjumătățire $=\ln 0\ds5/\ln(\text{persistență})$ zile.
Apoi, cu parametrii fixați, aplicați filtrul pe fereastra de test și comparați prognozele de varianță pentru o zi înainte, $\hat\sigma^2_t$, cu substitutul $\varepsilon_t^2$ prin funcția de pierdere de cvasi-verosimilitate
$\text{QLIKE}=\frac1n\sum_t\big(\ln\hat\sigma^2_t+\varepsilon_t^2/\hat\sigma^2_t\big)$ (mai mic înseamnă mai bine) \cite{Patton2011}.
La final traduceți în bani prognoza pentru o zi înainte $\hat\sigma_T$ a ultimei zile de test $T$: creșterea aproximativă a prețului de bază într-o zi, la cuantila de 99\%, $\bar P_{T-1}\,[\exp(z_{0\ds99}\hat\sigma_T/100)-1]$, unde $\bar P_{T-1}$ este prețul de bază realizat în ziua precedentă, iar $z_{0\ds99}=t^{-1}_\nu(0\ds99)\sqrt{(\nu-2)/\nu}$ este cuantila de 99\% a repartiției $t_\nu$ standardizate.}
\out{\T{table T7 (rows = 3 models; columns $\omega,\alpha,\gamma,\beta,\nu$, log-likelihood, AIC, BIC (Bayesian information criterion), persistence, half-life, $p$-value of $\gamma$, QLIKE on test); figure of $|\varepsilon_t|$ and the three $\hat\sigma_t$ on the test window; the 99\% move in EUR/MWh.}%
{tabelul T7 (rînduri = 3 modele; coloane $\omega,\alpha,\gamma,\beta,\nu$, log-verosimilitate, AIC, BIC (criteriul informațional bayesian), persistență, timp de înjumătățire, valoarea $p$ pentru $\gamma$, QLIKE pe test); figura cu $|\varepsilon_t|$ și cele trei $\hat\sigma_t$ pe fereastra de test; creșterea la cuantila de 99\%, în EUR/MWh.}}

\task{8}{\T{The memo: decision, reconciliation, your company}{Memoriul: decizie, reconciliere, compania dumneavoastră}}{22}
\T{\textbf{Page 1 (recommendation, at most 450 words):}}{\textbf{Pagina 1 (recomandarea, cel mult 450 de cuvinte):}}
\begin{enumerate}[label=(\alph*)]
\item \T{\textbf{Decision} in the first two lines: which model the desk should use for the focus block, from which date, and what it replaces.}%
{\textbf{Decizia} în primele două rînduri: ce model să folosească biroul pentru blocul de interes, de la ce dată și ce înlocuiește.}
\item \T{\textbf{Evidence} (three bullets, each with a number from your notebook): accuracy gain vs naive (rMAE, EUR/MWh), DM $p$-value, the 90\% interval width and its coverage.}%
{\textbf{Dovezile} (trei enumerări, fiecare cu un număr din notebook): cîștigul de acuratețe față de modelul naiv (rMAE, EUR/MWh), valoarea $p$ DM, lățimea intervalului de 90\% și acoperirea lui.}
\item \T{\textbf{Risk}: what the GARCH result says about the current regime, and your buffer recommendation in EUR/MWh.}%
{\textbf{Riscul}: ce spune rezultatul GARCH despre regimul actual și recomandarea privind rezerva, în EUR/MWh.}
\item \T{\textbf{Limits}: the load/renewables proxy, one period where your model failed, and what you would add next.}%
{\textbf{Limitele}: substitutul pentru consum/regenerabile, o perioadă în care modelul a eșuat și ce ați adăuga în continuare.}
\item \T{\textbf{How this applies in my company} (at most 120 words). Public information only: never use or disclose confidential data. Example: ``our supplier buys about X\% of its portfolio on PZU in the evening peak''
is acceptable only if X is public; otherwise describe the mechanism without numbers.}%
{\textbf{Cum se aplică în compania mea} (cel mult 120 de cuvinte). Numai informații publice: nu folosiți și nu divulgați date confidențiale. Exemplu: „furnizorul nostru cumpără circa X\% din portofoliu pe PZU la vîrful de seară''
este acceptabil numai dacă X este public; altfel descrieți mecanismul fără cifre.}
\end{enumerate}
\T{\textbf{Page 2 (appendix):} your parameter block; excerpts of T5, T6, T7; and the \textbf{hand-reconciliation table}:
for each of your two \texttt{check\_days} at \texttt{check\_hour}, write by hand the actual price, the Naive, AW-ARX and LEAR forecasts, each absolute error,
the LEAR+CP 5\% and 95\% quantiles, whether the price was inside, and the pinball score at $\tau=0\ds95$, with the arithmetic shown (e.g.\ $(0-0\ds95)(278\ds9-394\ds9)=110\ds2$).}%
{\textbf{Pagina 2 (anexa):} blocul de parametri; extrase din T5, T6, T7; și \textbf{tabelul de reconciliere manuală}:
pentru fiecare dintre cele două \texttt{check\_days} la \texttt{check\_hour}, scrieți de mînă prețul realizat, prognozele Naive, AW-ARX și LEAR, fiecare eroare absolută,
cuantilele de 5\% și 95\% LEAR+CP, dacă prețul a fost în interval și scorul pinball la $\tau=0\ds95$, cu calculul arătat (de exemplu $(0-0\ds95)(278\ds9-394\ds9)=110\ds2$).}

\section{\T{Exact conventions (answers to the questions students usually ask)}{Convenții exacte (răspunsuri la întrebările frecvente ale studenților)}}
\small
\begin{enumerate}[label=\textbf{C\arabic*}]
\item \T{\textbf{Names in \texttt{F}} (a Python \texttt{dict} of \texttt{DataFrame}s, rows \texttt{fdays}, columns \texttt{h00}..\texttt{h23}): \texttt{"Naive (Weron)"}, \texttt{"Seasonal naive d-7"}, \texttt{"ARX(56)"}, \dots, \texttt{"ARX(728)"}, \texttt{"AW-ARX(6 windows)"}, \texttt{"LEAR(<lear\_window>)"}. Columns of T5 per year: \texttt{MAE\_2025}, \texttt{RMSE\_2025}, \texttt{rMAE\_2025}.}%
{\textbf{Numele din \texttt{F}} (un \texttt{dict} Python de \texttt{DataFrame}-uri, rînduri \texttt{fdays}, coloane \texttt{h00}..\texttt{h23}): \texttt{"Naive (Weron)"}, \texttt{"Seasonal naive d-7"}, \texttt{"ARX(56)"}, \dots, \texttt{"ARX(728)"}, \texttt{"AW-ARX(6 windows)"}, \texttt{"LEAR(<lear\_window>)"}. Coloanele T5 pe ani: \texttt{MAE\_2025}, \texttt{RMSE\_2025}, \texttt{rMAE\_2025}.}
\item \T{\textbf{VST}: $m$ and $s$ are computed once per forecast day and window, \emph{pooled over all 24 hours} of the price rows $d-T-7,\dots,d-1$ (targets plus lag rows; hence 63 days in the worked example for $T=56$), and used for every hour.
\textbf{ARX exogenous}: $L$ and $R$ are z-scored with one mean and one standard deviation pooled over the same rows and hours, but only up to row $d-2$ (no $\asinh$).
\textbf{LEAR exogenous}: $L$ and $R$ get their own $\asinh$ transform (own median/MAD, same rows up to $d-2$), then every one of the 247 columns, dummies included, is z-scored with calibration means and standard deviations. For $h=23$ the ARX regressor $Y_{d-1,23}$ appears twice: drop the duplicate (11 regressors).}%
{\textbf{VST}: $m$ și $s$ se calculează o dată pentru fiecare zi de prognoză și fereastră, \emph{comun pentru toate cele 24 de ore} ale rîndurilor de preț $d-T-7,\dots,d-1$ (ținte plus rîndurile decalajelor; de aceea 63 de zile în exemplul rezolvat pentru $T=56$) și se folosesc pentru fiecare oră.
\textbf{Variabilele exogene ARX}: $L$ și $R$ se standardizează cu o singură medie și o singură abatere standard, comune pentru aceleași rînduri și ore, dar numai pînă la rîndul $d-2$ (fără $\asinh$).
\textbf{Variabilele exogene LEAR}: $L$ și $R$ primesc propria transformare $\asinh$ (mediană/MAD proprii, aceleași rînduri pînă în $d-2$), apoi fiecare dintre cele 247 de coloane, inclusiv variabilele zilei, se standardizează cu mediile și abaterile standard din calibrare. Pentru $h=23$, regresorul ARX $Y_{d-1,23}$ apare de două ori: eliminați dublura (11 regresori).}
\item \T{\textbf{Years and quarters}: calendar years and calendar quarters intersected with your window (partial periods allowed; report the number of days). rMAE always divides by the Naive MAE on the same subset (same days, same hours).}%
{\textbf{Ani și trimestre}: ani și trimestre calendaristice intersectate cu fereastra (perioadele parțiale sînt permise; raportați numărul de zile). rMAE se împarte întotdeauna la MAE naiv pe aceeași submulțime (aceleași zile, aceleași ore).}
\item \T{\textbf{DM matrix}: the cell in row A, column B holds $p=1-\Phi(\text{DM})$ with $\Delta_d=$ mean $|e^A|-$ mean $|e^B|$; the diagonal is empty.}%
{\textbf{Matricea DM}: celula din rîndul A, coloana B conține $p=1-\Phi(\text{DM})$ cu $\Delta_d=$ media $|e^A|-$ media $|e^B|$; diagonala este goală.}
\item \T{\textbf{Kupiec on the focus block}: $n$ = test days $\times$ number of focus hours; $x$ = (day, hour) pairs outside $[\hat q^{0\ds05},\hat q^{0\ds95}]$ on either side; ``inside'' includes the bounds; if $x=0$ use $0\cdot\ln 0=0$. The pairs within a day are not independent: mention it.
\textbf{Best combination} (for the figure) = lowest APS on the test window; any test week. Report CRPS $\approx2\cdot$APS as an extra column.}%
{\textbf{Kupiec pe blocul de interes}: $n$ = zile de test $\times$ numărul orelor de interes; $x$ = perechile (zi, oră) în afara lui $[\hat q^{0\ds05},\hat q^{0\ds95}]$, de oricare parte; „în interval'' include capetele; dacă $x=0$ folosiți $0\cdot\ln 0=0$. Perechile din aceeași zi nu sînt independente: menționați acest lucru.
\textbf{Cea mai bună combinație} (pentru figură) = APS minim pe fereastra de test; orice săptămînă de test. Raportați CRPS $\approx2\cdot$APS ca o coloană suplimentară.}
\item \T{\textbf{GARCH in \texttt{arch}}: \texttt{garch\_years} are calendar years (\texttt{test\_start - pd.DateOffset(years=g)}). Subtract the estimation-sample weekday means from all $r_t$ (estimation and test).
Fit \texttt{arch\_model(r, mean="Constant", vol="GARCH", p=1, o=0 or 1, q=1, dist="t")} on the series up to \texttt{test\_end} with \texttt{.fit(last\_obs=test\_start)} (parameters use only the estimation sample), then \texttt{.forecast(horizon=1, start=<last estimation day>, reindex=False)}: the forecast made at day $t$ is the variance of day $t+1$. For EGARCH use \texttt{vol="EGARCH"}. The $p$-value of $\gamma$ is ``n/a'' for GARCH(1,1).}%
{\textbf{GARCH în \texttt{arch}}: \texttt{garch\_years} sînt ani calendaristici (\texttt{test\_start - pd.DateOffset(years=g)}). Scădeți mediile pe zile ale săptămînii din eșantionul de estimare din toate valorile $r_t$ (estimare și test).
Estimați \texttt{arch\_model(r, mean="Constant", vol="GARCH", p=1, o=0 sau 1, q=1, dist="t")} pe seria pînă la \texttt{test\_end} cu \texttt{.fit(last\_obs=test\_start)} (parametrii folosesc doar eșantionul de estimare), apoi \texttt{.forecast(horizon=1, start=<ultima zi de estimare>, reindex=False)}: prognoza făcută în ziua $t$ este varianța zilei $t+1$. Pentru EGARCH folosiți \texttt{vol="EGARCH"}. Valoarea $p$ pentru $\gamma$ este „n/a'' la GARCH(1,1).}
\item \T{\textbf{Leakage check}: run \texttt{ea.leakage\_check} at least for Naive, ARX(56) and LEAR on the first test day.
\textbf{Where to work}: keep your notebook in \texttt{repo/assignments/A1/} next to the starter (it imports \texttt{emqa\_assign} from the parent folder).
\textbf{\texttt{MODE="latest"}}: only a clean run is required; tables may contain empty cells (e.g.\ one quarter only).
\textbf{Hand reconciliation}: typed arithmetic is fine; the pinball score is that of the LEAR+CP quantile $\hat q^{0\ds95}$.}%
{\textbf{Verificarea scurgerii}: rulați \texttt{ea.leakage\_check} cel puțin pentru Naive, ARX(56) și LEAR în prima zi de test.
\textbf{Unde lucrați}: păstrați notebook-ul în \texttt{repo/assignments/A1/}, lîngă cel de pornire (importă \texttt{emqa\_assign} din dosarul părinte).
\textbf{\texttt{MODE="latest"}}: se cere doar o rulare fără erori; tabelele pot avea celule goale (de exemplu un singur trimestru).
\textbf{Reconcilierea manuală}: calculele tastate sînt acceptate; scorul pinball este cel al cuantilei LEAR+CP $\hat q^{0\ds95}$.}
\end{enumerate}
\normalsize

\section{\T{Worked example (real data): the rolling window and an HS/CP interval}{Exemplu rezolvat (date reale): fereastra glisantă și un interval HS/CP}}
\T{\textbf{Forecast day} $d=$ Tuesday 4 March 2025, hour \texttt{h19} (19:00--20:00 CET), information cut-off 3 March 2025, 12:00 CET.}%
{\textbf{Ziua de prognoză} $d=$ marți, 4 martie 2025, ora \texttt{h19} (19:00--20:00 CET), momentul-limită al informației: 3 martie 2025, ora 12:00 CET.}

\T{\textbf{(1) Calibration window for ARX(56).}}{\textbf{(1) Fereastra de calibrare pentru ARX(56).}}
\begin{center}\small
\begin{tabular}{l l l}
\toprule
\T{Item}{Element} & \T{Dates}{Date} & \T{Why}{De ce}\\
\midrule
\T{target days (56)}{zile-țintă (56)} & 2025-01-07 \dots\ 2025-03-03 & $d-56,\dots,d-1$\\
\T{oldest price row used}{cel mai vechi rînd de preț folosit} & 2024-12-31 & \T{lag 7 of the first target}{decalajul 7 al primei ținte}\\
$m,s$ \T{of the VST}{ale VST} & \T{prices 2024-12-31 \dots\ 2025-03-03}{prețurile 2024-12-31 \dots\ 2025-03-03} & \T{never the test day}{niciodată ziua de test}\\
\T{$L$, $R$ (proxies, scaling)}{$L$, $R$ (substitute, standardizare)} & \T{rows up to 2025-03-02}{rîndurile pînă la 2025-03-02} & \T{$d-1$ not yet metered}{$d-1$ nu este încă măsurată}\\
\T{next step}{pasul următor} & \T{targets 2025-01-08 \dots\ 2025-03-04}{ținte 2025-01-08 \dots\ 2025-03-04} & \T{forecast 5 March (roll one day)}{prognoza pentru 5 martie}\\
\bottomrule
\end{tabular}
\end{center}
\T{\textbf{(2) Regressors of $d$ for hour h19 (before the VST), from the course data.}}{\textbf{(2) Regresorii zilei $d$ pentru ora h19 (înainte de VST), din datele cursului.}}
\begin{center}\small\setlength{\tabcolsep}{3.5pt}
\begin{tabular}{cccccccccc}
\toprule
$P_{d-1,19}$ & $P_{d-2,19}$ & $P_{d-7,19}$ & $P_{d-1,23}$ & $P^{\min}_{d-1}$ & $P^{\max}_{d-1}$ & $L_{d-2,19}$ & $R_{d-2,19}$ & $D^{\text{Mon}},D^{\text{Sat}},D^{\text{Sun}}$ & \T{actual}{realizat} $P_{d,19}$\\
\midrule
198\ds18 & 178\ds86 & 438\ds63 & 151\ds23 & 46\ds34 & 286\ds21 & 7\,119 MW & 843 MW & 0/0/0 & 231\ds65\\
\bottomrule
\end{tabular}
\end{center}
\T{The naive forecast is $P_{d-1,19}=198\ds18$ (Tuesday uses $d-1$), error $231\ds65-198\ds18=33\ds47$ EUR/MWh.}%
{Prognoza naivă este $P_{d-1,19}=198\ds18$ (marțea folosește $d-1$), eroarea $231\ds65-198\ds18=33\ds47$ EUR/MWh.}

\T{\textbf{(3) An 80\% interval from the last $W=10$ naive errors at h19} (22 Feb -- 3 Mar 2025; $W=10$ only to keep the example short):}%
{\textbf{(3) Un interval de 80\% din ultimele $W=10$ erori naive la h19} (22 feb.\ -- 3 mar.\ 2025; $W=10$ numai ca exemplul să fie scurt):}
\begin{center}\small
\begin{tabular}{l r r r r r r r r r r}
\toprule
\T{sorted errors}{erorile sortate} & $-246\ds5$ & $-159\ds4$ & $-66\ds0$ & $-16\ds8$ & $-13\ds3$ & $55\ds1$ & $64\ds3$ & $66\ds0$ & $81\ds1$ & $140\ds7$\\
\T{sorted $|e|$}{$|e|$ sortate} & $13\ds3$ & $16\ds8$ & $55\ds1$ & $64\ds3$ & $66\ds0$ & $66\ds0$ & $81\ds1$ & $140\ds7$ & $159\ds4$ & $246\ds5$\\
\bottomrule
\end{tabular}
\end{center}
\T{HS: $Q_{0\ds10}(e)=-168\ds1$ and $Q_{0\ds90}(e)=87\ds0$ (linear interpolation), so the interval is $[198\ds2-168\ds1,\;198\ds2+87\ds0]=[30\ds1,\;285\ds2]$.
CP: $Q_{0\ds80}(|e|)=144\ds4$, so $[198\ds2-144\ds4,\;198\ds2+144\ds4]=[53\ds7,\;342\ds6]$. The actual 231\ds65 lies inside both.
Pinball at $\tau=0\ds90$ for HS: $\hat q=285\ds2>y=231\ds65$, so $\text{PS}=(1-0\ds9)(285\ds2-231\ds65)=5\ds36$.}%
{HS: $Q_{0\ds10}(e)=-168\ds1$ și $Q_{0\ds90}(e)=87\ds0$ (interpolare liniară), deci intervalul este $[198\ds2-168\ds1,\;198\ds2+87\ds0]=[30\ds1,\;285\ds2]$.
CP: $Q_{0\ds80}(|e|)=144\ds4$, deci $[198\ds2-144\ds4,\;198\ds2+144\ds4]=[53\ds7,\;342\ds6]$. Valoarea realizată 231\ds65 este în ambele.
Scorul pinball la $\tau=0\ds90$ pentru HS: $\hat q=285\ds2>y=231\ds65$, deci $\text{PS}=(1-0\ds9)(285\ds2-231\ds65)=5\ds36$.}

\T{\textbf{(4) DM in one line.} If on a test day the mean $|e|$ over your focus hours is 30 for Naive (A) and 22 for LEAR (B), then $\Delta_d=+8$:
LEAR was better that day. The test asks whether the mean of these daily differences is significantly above zero.}%
{\textbf{(4) DM într-un rînd.} Dacă într-o zi de test media lui $|e|$ pe orele de interes este 30 pentru Naive (A) și 22 pentru LEAR (B), atunci $\Delta_d=+8$:
LEAR a fost mai bun în acea zi. Testul verifică dacă media acestor diferențe zilnice este semnificativ peste zero.}

\section{\T{Rules: reproducibility and use of AI tools}{Reguli: reproductibilitate și folosirea instrumentelor de inteligență artificială}}
\begin{enumerate}[label=\textbf{R\arabic*}]
\item \T{\textbf{Runs end-to-end}: \emph{Restart \& Run All} works on the course data with the packages of the course environment (\texttt{numpy}, \texttt{pandas}, \texttt{scikit-learn}, \texttt{statsmodels}, \texttt{arch}, \texttt{matplotlib}).}%
{\textbf{Rulează complet}: \emph{Restart \& Run All} funcționează pe datele cursului cu pachetele mediului cursului (\texttt{numpy}, \texttt{pandas}, \texttt{scikit-learn}, \texttt{statsmodels}, \texttt{arch}, \texttt{matplotlib}).}
\item \T{\textbf{No hard-coded dates or numbers}: dates come from \texttt{p}. The instructor re-runs every notebook with \texttt{MODE = "latest"} on data downloaded after the deadline (an unseen hold-out); it must run.}%
{\textbf{Fără date sau numere introduse direct în cod (\emph{hard-coded})}: datele calendaristice vin din \texttt{p}. Profesorul rulează din nou fiecare notebook cu \texttt{MODE = "latest"} pe date descărcate după termen (un set de verificare nevăzut); trebuie să ruleze.}
\item \T{\textbf{Data version}: the notebook prints \texttt{ea.data\_version()} (file hashes).}{\textbf{Versiunea datelor}: notebook-ul afișează \texttt{ea.data\_version()} (amprentele fișierelor).}
\item \T{\textbf{No look-ahead}: \texttt{ea.leakage\_check} passes for every model you report.}{\textbf{Fără informații din viitor}: \texttt{ea.leakage\_check} trece pentru fiecare model raportat.}
\item \T{\textbf{Every number in the memo is in \texttt{MY\_NUMBERS}}. A number that is not, or that differs from the notebook, costs 2 points (up to 10).}%
{\textbf{Fiecare număr din memoriu se află în \texttt{MY\_NUMBERS}}. Un număr care nu este acolo sau care diferă de notebook costă 2 puncte (pînă la 10).}
\item \T{\textbf{AI tools} (ChatGPT, Claude, Copilot, \dots) are allowed for code, with disclosure: add a cell ``AI use'' that names the tool and what it did.
They are not allowed to write the memo's numbers or conclusions. You must be able to explain every line at the oral defence.}%
{\textbf{Instrumentele de inteligență artificială} (ChatGPT, Claude, Copilot, \dots) sînt permise pentru cod, cu declarare: adăugați o celulă „Folosirea AI'' care numește instrumentul și ce a făcut.
Nu au voie să scrie numerele sau concluziile memoriului. Trebuie să puteți explica fiecare linie la susținerea orală.}
\end{enumerate}
\T{If the notebook does not run (R1), Tasks 1--7 are capped at 50\%. If you cannot explain a result at the oral defence, the instructor may set the points of that task to 0.}%
{Dacă notebook-ul nu rulează (R1), sarcinile 1--7 sînt plafonate la 50\%. Dacă nu puteți explica un rezultat la susținerea orală, profesorul poate acorda 0 puncte pentru acea sarcină.}

\section{\T{Common pitfalls}{Greșeli frecvente}}
\begin{warnbox}\small
\begin{itemize}
\item \T{\textbf{Look-ahead leakage}: using $L_{d,h}$ or $R_{d,h}$ of the forecast day; fitting the VST, the scaler or the LASSO on data that include day $d$; choosing a model on the test window and then reporting its test accuracy as if it were out-of-sample.}%
{\textbf{Scurgerea de informație din viitor}: folosirea lui $L_{d,h}$ sau $R_{d,h}$ din ziua prognozată; estimarea VST, a standardizării sau a LASSO pe date care includ ziua $d$; alegerea modelului pe fereastra de test și raportarea acurateței pe același test ca și cum ar fi în afara eșantionului.}
\item \T{\textbf{Logarithms of prices} or log-returns: RO had 196 negative hours in 2025 alone. Use $\asinh$.}{\textbf{Logaritmarea prețurilor} sau randamentele logaritmice: România a avut 196 de ore cu prețuri negative numai în 2025. Folosiți $\asinh$.}
\item \T{\textbf{MAPE} with prices near zero; deleting negative prices to make the errors look smaller.}{\textbf{MAPE} cu prețuri apropiate de zero; ștergerea prețurilor negative ca erorile să pară mai mici.}
\item \T{\textbf{Mixing 15-minute and hourly data}, or UTC and CET hours (a one-hour shift destroys the hourly models).}{\textbf{Amestecarea datelor de 15 minute cu cele orare} sau a orelor UTC cu cele CET (un decalaj de o oră strică modelele orare).}
\item \T{\textbf{In-sample fit reported as forecast accuracy}; $R^2$ instead of MAE/RMSE on the test window.}{\textbf{Ajustarea în eșantion raportată drept acuratețe a prognozei}; $R^2$ în loc de MAE/RMSE pe fereastra de test.}
\item \T{\textbf{Reading DM backwards}: our $p$-value is for ``column better than row''.}{\textbf{Citirea inversă a testului DM}: valoarea noastră $p$ este pentru „coloana mai bună decît rîndul''.}
\item \T{\textbf{Generic text}: ``electricity prices are volatile'' earns nothing. Every claim needs a number from your window.}{\textbf{Text generic}: „prețurile energiei sînt volatile'' nu aduce puncte. Fiecare afirmație are nevoie de un număr din fereastra dumneavoastră.}
\end{itemize}
\end{warnbox}

\section{\T{Before you submit: checklist}{Înainte de predare: lista de verificare}}
\begin{itemize}[label=$\square$]
\item \T{Parameter block printed in the notebook and copied into the memo appendix.}{Blocul de parametri afișat în notebook și copiat în anexa memoriului.}
\item \T{\emph{Restart \& Run All} works; run time noted; \texttt{MODE = "latest"} also runs.}{\emph{Restart \& Run All} funcționează; timpul de rulare notat; și \texttt{MODE = "latest"} rulează.}
\item \T{\texttt{ea.data\_version()} printed; \texttt{ea.leakage\_check} passed for Naive, ARX(56), LEAR.}{\texttt{ea.data\_version()} afișat; testul \texttt{ea.leakage\_check} reușit pentru Naive, ARX(56), LEAR.}
\item \T{Tables T1, T5, T6, T7, two DM matrices, three figures (axes with units, legend below).}{Tabelele T1, T5, T6, T7, cele două matrice DM, trei figuri (axe cu unități, legenda dedesubt).}
\item \T{Memo $\le$ 2 pages; decision in the first two lines; every number traceable to \texttt{MY\_NUMBERS}.}{Memoriul $\le$ 2 pagini; decizia în primele două rînduri; fiecare număr se regăsește în \texttt{MY\_NUMBERS}.}
\item \T{Hand-reconciliation table with arithmetic; ``AI use'' cell; no confidential company data.}{Tabelul de reconciliere manuală cu calcule; celula „Folosirea AI''; fără date confidențiale ale companiei.}
\end{itemize}

\section{\T{Frequently asked questions}{Întrebări frecvente}}
\small
\begin{description}[style=nextline, leftmargin=0pt, itemsep=2pt]
\item[\T{Which time zone?}{Ce fus orar?}] \T{CET/CEST everywhere (the files are already converted). Bucharest time = CET + 1 h. Do not convert again.}{CET/CEST peste tot (fișierele sînt deja convertite). Ora Bucureștiului = CET + 1 h. Nu le mai convertiți.}
\item[\T{Can I forecast the 96 quarter-hours?}{Pot prognoza cele 96 de sferturi de oră?}] \T{Not required. The models and the grading use the hourly matrix. A 15-minute analysis is welcome as an extra figure but earns no extra points.}{Nu este obligatoriu. Modelele și notarea folosesc matricea orară. O analiză pe 15 minute este binevenită ca figură suplimentară, dar nu aduce puncte în plus.}
\item[\T{What about negative prices and the price floor?}{Ce fac cu prețurile negative și cu limita minimă?}] \T{Keep them; they are real. The harmonised minimum clearing price of the day-ahead coupling (SDAC) is $-500$ EUR/MWh and became $-600$ EUR/MWh for delivery from 29 May 2026 (OPCOM notice of 7 May 2026).}{Le păstrați; sînt reale. Prețul minim armonizat de echilibru (\emph{clearing}) al cuplării pentru ziua următoare (SDAC) este $-500$ EUR/MWh și a devenit $-600$ EUR/MWh pentru livrările de la 29 mai 2026 (anunțul OPCOM din 7 mai 2026).}
\item[\T{Are there missing days?}{Există zile lipsă?}] \T{Not in \texttt{ea.panel()} from 2022 on (checked). If your code finds a gap, you probably mixed time zones. Days before 2019 have gaps; you do not need them.}{Nu în \texttt{ea.panel()} începînd cu 2022 (verificat). Dacă codul găsește un gol, probabil ați amestecat fusurile orare. Zilele dinainte de 2019 au goluri; nu aveți nevoie de ele.}
\item[\T{Why the burn-in?}{De ce perioada de rodaj?}] \T{HS and CP need $W$ past forecast errors. Without the 112 burn-in days your first test days would have no interval. Accuracy is measured on the test window only.}{HS și CP au nevoie de $W$ erori de prognoză anterioare. Fără cele 112 zile de rodaj, primele zile de test nu ar avea interval. Acuratețea se măsoară numai pe fereastra de test.}
\item[\T{LEAR is slow. What can I do?}{LEAR merge încet. Ce pot face?}] \T{About 1 minute on 8 cores with \texttt{joblib.Parallel} over days; 10--20 minutes on one core. Weekly recalibration is allowed if declared.}{Circa 1 minut pe 8 nuclee cu \texttt{joblib.Parallel} pe zile; 10--20 de minute pe un singur nucleu. Recalibrarea săptămînală este permisă dacă o declarați.}
\item[\T{Can I add XGBoost, Prophet or a neural network?}{Pot adăuga XGBoost, Prophet sau o rețea neurală?}] \T{Yes, in addition to (not instead of) the required models, with the same rolling scheme. No extra points in Assignment 1; machine learning is graded in Assignment 2.}{Da, în plus față de modelele cerute (nu în locul lor), cu aceeași schemă glisantă. Nu aduce puncte în plus la tema 1; învățarea automată este notată la tema 2.}
\item[\T{My DM $p$-value is 1\ds 000. Is that an error?}{Valoarea $p$ DM este 1\ds 000. Este o eroare?}] \T{No. It means the column model is clearly \emph{worse} than the row model; read the symmetric cell.}{Nu. Înseamnă că modelul de pe coloană este clar \emph{mai slab} decît cel de pe rînd; citiți celula simetrică.}
\item[\T{Which quantile convention?}{Ce convenție pentru cuantile?}] \T{\texttt{numpy.quantile} default (linear). Small differences from other conventions are not penalised.}{Cea implicită din \texttt{numpy.quantile} (liniară). Diferențele mici față de alte convenții nu se depunctează.}
\item[\T{Why $\asinh$ for GARCH and not percentage returns?}{De ce $\asinh$ pentru GARCH și nu randamente procentuale?}] \T{Percentage changes explode when $\bar P_{t-1}$ is near zero. For prices well above 1 EUR/MWh, $\asinh(P)\approx\ln(2P)$, so $r_t$ behaves like a log-return.}{Variațiile procentuale explodează cînd $\bar P_{t-1}$ este aproape de zero. Pentru prețuri mult peste 1 EUR/MWh, $\asinh(P)\approx\ln(2P)$, deci $r_t$ se comportă ca un randament logaritmic.}
\item[\T{My GJR $\gamma$ is positive, but the lecture said ``inverse leverage''.}{$\gamma$ din GJR este pozitiv, dar la curs s-a vorbit de „levier invers''.}] \T{Report what your data say and discuss why (e.g.\ daily base prices after weekday adjustment vs.\ hourly spikes). There is no ``right sign'' for the grade.}{Raportați ce arată datele și discutați de ce (de exemplu prețuri zilnice de bază după ajustarea pentru ziua săptămînii vs.\ salturi orare). Nu există un „semn corect'' pentru notă.}
\item[\T{Can I use R or Excel?}{Pot folosi R sau Excel?}] \T{The graded notebook must be Python (TODO: instructor may allow R).}{Notebook-ul notat trebuie să fie Python (TODO: profesorul poate permite R).}
\item[\T{Team work?}{Lucru în echipă?}] \T{If allowed (see the box on page 1), use \texttt{-{}-id "ID1+ID2"}; both members defend orally and must explain any part.}{Dacă este permis (vezi caseta de pe pagina 1), folosiți \texttt{-{}-id "ID1+ID2"}; ambii membri susțin oral și trebuie să explice orice parte.}
\end{description}
\normalsize

\section{\T{Grading rubric (100 points)}{Grila de notare (100 de puncte)}}
\small\setlength{\tabcolsep}{4pt}
\begin{longtable}{>{\raggedright\arraybackslash}p{2.4cm} >{\centering\arraybackslash}p{0.8cm} >{\raggedright\arraybackslash}p{4.4cm} >{\raggedright\arraybackslash}p{3.8cm} >{\raggedright\arraybackslash}p{3.6cm}}
\toprule
\T{Task}{Sarcina} & \T{Pts}{Pct.} & \T{Excellent (90--100\%)}{Excelent (90--100\%)} & \T{Adequate (50--70\%)}{Adecvat (50--70\%)} & \T{Poor (0--30\%)}{Slab (0--30\%)}\\
\midrule\endhead
\T{1 Test window}{1 Fereastra de test} & 8 & \T{T1 complete incl.\ negative MTUs; clean figure; events named and dated from the data}{T1 complet, inclusiv intervalele MTU negative; figură curată; evenimente numite și datate din date} & \T{table complete, events generic}{tabel complet, evenimente generice} & \T{wrong window or missing table}{fereastră greșită sau tabel lipsă}\\
\T{2 Naive}{2 Naive} & 6 & \T{both rules exact (weekday logic), leakage check passed}{ambele reguli exacte (logica zilelor), verificarea scurgerii trecută} & \T{one rule wrong on some weekdays}{o regulă greșită în unele zile} & \T{missing or uses day $d$}{lipsește sau folosește ziua $d$}\\
\T{3 ARX + averaging}{3 ARX + mediere} & 12 & \T{exact regressors, VST fitted in-window, 6 rolling windows, daily recalibration, AW}{regresori exacți, VST estimat în fereastră, 6 ferestre glisante, recalibrare zilnică, AW} & \T{fixed window or some regressors missing}{fereastră fixă sau lipsesc unii regresori} & \T{in-sample fit only, or leakage}{doar ajustare în eșantion sau scurgere de informație}\\
\T{4 LEAR}{4 LEAR} & 10 & \T{247 regressors, AIC-selected $\lambda$ per hour, rolling, sparsity reported}{247 de regresori, $\lambda$ ales prin AIC pentru fiecare oră, glisant, numărul de coeficienți raportat} & \T{fewer regressors or fixed $\lambda$}{mai puțini regresori sau $\lambda$ fix} & \T{not rolling / not run}{nu este glisant / nu rulează}\\
\T{5 Accuracy, DM}{5 Acuratețe, DM} & 16 & \T{all metrics by year, quarter and focus block; HAC DM both matrices; correct reading; significance discussed}{toate măsurile pe ani, trimestre și bloc; ambele matrice DM HAC; citire corectă; semnificația discutată} & \T{metrics right, DM missing one matrix or misread once}{măsuri corecte, o matrice DM lipsă sau citită greșit o dată} & \T{metrics on burn-in, MAPE, no DM}{măsuri pe rodaj, MAPE, fără DM}\\
\T{6 Probabilistic}{6 Probabilistic} & 14 & \T{6 combinations, correct HS/CP, APS, coverage, Kupiec, figure; calibration vs sharpness discussed}{6 combinații, HS/CP corecte, APS, acoperire, Kupiec, figură; calibrarea vs.\ precizia discutate} & \T{HS or CP only, no Kupiec}{numai HS sau CP, fără Kupiec} & \T{intervals from in-sample residuals or future errors}{intervale din reziduuri în eșantion sau din erori viitoare}\\
\T{7 Volatility}{7 Volatilitate} & 12 & \T{3 models with $t$ errors, $\gamma$ tested and interpreted, persistence, QLIKE out-of-sample, 99\% move in EUR}{3 modele cu erori $t$, $\gamma$ testat și interpretat, persistență, QLIKE în afara eșantionului, creșterea la cuantila de 99\% în EUR} & \T{in-sample only or no money translation}{numai în eșantion sau fără traducere în bani} & \T{log-returns of prices; no estimation window}{randamente logaritmice ale prețurilor; fără fereastră de estimare}\\
\T{8a Decision}{8a Decizia} & 10 & \T{clear decision first; 3 numbers; limits and failure period honest}{decizie clară la început; 3 numere; limite și perioada de eșec onest prezentate} & \T{decision vague or numbers without meaning}{decizie vagă sau numere fără semnificație} & \T{generic text, no decision}{text generic, fără decizie}\\
\T{8b Reconciliation}{8b Reconciliere} & 6 & \T{both days by hand, arithmetic shown, matches notebook}{ambele zile de mînă, calcule arătate, coincid cu notebook-ul} & \T{one day or arithmetic missing}{o zi sau calcule lipsă} & \T{missing or inconsistent}{lipsă sau inconsistent}\\
\T{8c My company}{8c Compania mea} & 6 & \T{concrete mechanism in the student's job, public data only}{mecanism concret în postul studentului, numai date publice} & \T{plausible but generic}{plauzibil, dar generic} & \T{absent, or confidential data}{lipsă sau date confidențiale}\\
\bottomrule
\end{longtable}
\normalsize
\T{Penalties and gates: R1 (does not run) caps Tasks 1--7 at 50\%; R5 costs 2 points per untraceable number (up to 10); wrong personal parameters = 0.
A short oral defence confirms authorship and may reduce, never increase, the points of a task.}%
{Penalizări și condiții: R1 (nu rulează) plafonează sarcinile 1--7 la 50\%; R5 costă 2 puncte pentru fiecare număr care nu se regăsește (pînă la 10); parametri personali greșiți = 0.
Susținerea orală scurtă confirmă paternitatea lucrării și poate reduce, niciodată crește, punctajul unei sarcini.}

\section{\T{Further reading (papers of the course instructor and co-authors)}{Lecturi suplimentare (articole ale titularului de curs și ale coautorilor)}}
\small
\begin{itemize}
\item A.-V.~Andrei, G.~Velev, F.-M.~Toma, D.T.~Pele, S.~Lessmann (2024). \emph{Energy Price Modelling: A Comparative Evaluation of four Generations of Forecasting Methods}, \href{https://arxiv.org/abs/2411.03372}{arXiv:2411.03372} --- \T{Tasks 2--5}{sarcinile 2--5}.
\item \emph{Explainable AI-Based Early Warning System for Electricity Price Spikes: Evidence from the Romanian Market}, PICBE 2026, \href{https://doi.org/10.2478/picbe-2026-0352}{doi:10.2478/picbe-2026-0352} --- \T{Tasks 1 and 7}{sarcinile 1 și 7}.
\item \emph{Evaluating Intervention Effectiveness during Extreme Price Events: Romania's 2024 Energy Crisis}, PICBE 2026, \href{https://doi.org/10.2478/picbe-2026-0013}{doi:10.2478/picbe-2026-0013} --- \T{Task 1}{sarcina 1}.
\end{itemize}
\normalsize
\begin{thebibliography}{99}\small
\bibitem{Diebold1995} F.X.~Diebold, R.S.~Mariano (1995). Comparing predictive accuracy. \emph{Journal of Business \& Economic Statistics} 13(3), 253--263, \href{https://doi.org/10.1080/07350015.1995.10524599}{doi:10.1080/07350015.1995.10524599}.
\bibitem{Gneiting2007} T.~Gneiting, A.E.~Raftery (2007). Strictly proper scoring rules, prediction, and estimation. \emph{Journal of the American Statistical Association} 102, 359--378, \href{https://doi.org/10.1198/016214506000001437}{doi:10.1198/016214506000001437}.
\bibitem{Hubicka2019} K.~Hubicka, G.~Marcjasz, R.~Weron (2019). A note on averaging day-ahead electricity price forecasts across calibration windows. \emph{IEEE Transactions on Sustainable Energy} 10(1), 321--323, \href{https://doi.org/10.1109/TSTE.2018.2869557}{doi:10.1109/TSTE.2018.2869557}.
\bibitem{Kupiec1995} P.~Kupiec (1995). Techniques for verifying the accuracy of risk measurement models. \emph{Journal of Derivatives} 3, 73--84, \href{https://doi.org/10.3905/jod.1995.407942}{doi:10.3905/jod.1995.407942}.
\bibitem{Lago2021} J.~Lago, G.~Marcjasz, B.~De Schutter, R.~Weron (2021). Forecasting day-ahead electricity prices: A review of state-of-the-art algorithms, best practices and an open-access benchmark. \emph{Applied Energy} 293, 116983 (\texttt{epftoolbox}), \href{https://doi.org/10.1016/j.apenergy.2021.116983}{doi:10.1016/j.apenergy.2021.116983}.
\bibitem{Lipiecki2024} A.~Lipiecki, B.~Uniejewski, R.~Weron (2024). Postprocessing of point predictions for probabilistic forecasting of day-ahead electricity prices: The benefits of using isotonic distributional regression. \emph{Energy Economics} 139, 107934, \href{https://doi.org/10.1016/j.eneco.2024.107934}{doi:10.1016/j.eneco.2024.107934}.
\bibitem{Nowotarski2018} J.~Nowotarski, R.~Weron (2018). Recent advances in electricity price forecasting: A review of probabilistic forecasting. \emph{Renewable and Sustainable Energy Reviews} 81, 1548--1568, \href{https://doi.org/10.1016/j.rser.2017.05.234}{doi:10.1016/j.rser.2017.05.234}.
\bibitem{Patton2011} A.J.~Patton (2011). Volatility forecast comparison using imperfect volatility proxies. \emph{Journal of Econometrics} 160, 246--256, \href{https://doi.org/10.1016/j.jeconom.2010.03.034}{doi:10.1016/j.jeconom.2010.03.034}.
\bibitem{Uniejewski2016} B.~Uniejewski, J.~Nowotarski, R.~Weron (2016). Automated variable selection and shrinkage for day-ahead electricity price forecasting. \emph{Energies} 9, 621, \href{https://doi.org/10.3390/en9080621}{doi:10.3390/en9080621}.
\bibitem{Uniejewski2018} B.~Uniejewski, R.~Weron, F.~Ziel (2018). Variance stabilizing transformations for electricity spot price forecasting. \emph{IEEE Transactions on Power Systems} 33(2), 2219--2229, \href{https://doi.org/10.1109/TPWRS.2017.2734563}{doi:10.1109/TPWRS.2017.2734563}.
\bibitem{Weron2014} R.~Weron (2014). Electricity price forecasting: A review of the state-of-the-art with a look into the future. \emph{International Journal of Forecasting} 30(4), 1030--1081, \href{https://doi.org/10.1016/j.ijforecast.2014.08.008}{doi:10.1016/j.ijforecast.2014.08.008}.
\bibitem{WeronIIF} R.~Weron (2025). \emph{Electricity Price Forecasting}, IIF Distinguished Lectures, parts I--II, ISEA 2025, \href{https://p.wz.pwr.edu.pl/~weron.rafal/Conf/IIF24}{\T{slides and recordings}{prezentări și înregistrări}}, and \emph{Machine learning for electricity price forecasting}, keynote, \href{https://www.ai4efin.ase.ro/aida-conference/}{AIDA 2026}, Bucharest.
\bibitem{Ziel2018} F.~Ziel, R.~Weron (2018). Day-ahead electricity price forecasting with high-dimensional structures: Univariate vs.\ multivariate modeling frameworks. \emph{Energy Economics} 70, 396--420, \href{https://doi.org/10.1016/j.eneco.2017.12.016}{doi:10.1016/j.eneco.2017.12.016}.
\end{thebibliography}
\end{document}
